\documentclass[11pt]{article}
\usepackage[margin=1in]{geometry}
\usepackage{amsmath,amssymb,amsthm,mathtools}
\usepackage[T1]{fontenc}
\usepackage{lmodern,microtype,enumitem,xurl}
\usepackage[colorlinks=true,linkcolor=blue,urlcolor=blue]{hyperref}
\hypersetup{pdftitle={Exponential sample complexity for an analytic activation arbitrarily close to tanh},pdfauthor={ChatGPT-6 Astra; Kaarel Hanni; Dmitry Vaintrob}}
\newtheorem{theorem}{Theorem}[section]
\newtheorem{lemma}[theorem]{Lemma}
\newtheorem{proposition}[theorem]{Proposition}
\theoremstyle{remark}\newtheorem{remark}[theorem]{Remark}
\newcommand{\R}{\mathbb R}
\newcommand{\E}{\mathbb E}
\newcommand{\norm}[1]{\left\lVert#1\right\rVert}
\newcommand{\poly}{\operatorname{poly}}
\newcommand{\Lip}{\operatorname{Lip}}
\newcommand{\dist}{\operatorname{dist}}
\title{Exponential sample complexity for an analytic activation\\
arbitrarily close to tanh}
\author{ChatGPT-6 Astra \and Kaarel H\"anni \and Dmitry Vaintrob\thanks{Authors are listed in alphabetical order.}}
\date{Research draft --- 25 October 2026}
\begin{document}
\maketitle
\begin{abstract}
For every fixed integer $r\ge1$ and every $\eta>0$, we construct one
computable odd, strictly increasing, real-analytic activation $\phi$
with range $(-1,1)$, global Lipschitz constant one, $\phi'(0)=1$,
strict concavity on the positive half-line, and
$\|\phi-\tanh\|_{C^r(\R)}<\eta$.
For this fixed activation, the original Gaussian Bayesian classifier
requires exponentially many examples in teacher depth on a uniform
Boolean cube. Teachers have depth $L_d=Dd+O(1)$, width and coefficient
bound $O(d)$, and constant global output margin. At every fixed
Gaussian exponent $\gamma\ge1$, every sample of size at most $\exp(cd)$
fails the posterior confidence criterion at every sufficiently large
width $N\ge\exp(C_\gamma d)$. The critical endpoint uses finite-width
entropy certificates and a parallel-circuit normal form for paths
near their scalar baseline. A fixed-width cold-limit statement also holds.
An appendix gives a second activation with every derivative globally
bounded and retains the $\exp(\Omega(L))$ lower bound.

The construction is computationally contrived. It encodes computable
hard truth tables in tiny analytic bumps, and decodes them by finite
differences that cancel the tanh background. No efficient evaluator
for the activation is asserted. This is a theorem about specially
constructed activations arbitrarily close to tanh, not about tanh
itself; it does not settle tanh, ReLU, or GELU.
\end{abstract}
\paragraph{Status.}
This conventional mathematical research proof has received independent
AI-agent reviews of the activation, decoder, capacity, energy, and posterior
arguments, including the critical finite-width target-mass and
canonical-state arguments. It is not Lean verified or human peer reviewed.

\section{Model and theorem}
Let $X_d=\{-1,1\}^d$ and $h^0(x)=(1,x)$. A common-width-$N$
network of weighted depth $L$ is
\[
 h^s(x)=\phi(W^s h^{s-1}(x))\quad(1\le s<L),\qquad
 z_W(x)=v^{\mathsf T}h^{L-1}(x).
\]
There are no separate hidden biases. Write $W^L=v^{\mathsf T}$
and $P_F(W)=\prod_{s=1}^L\|W^s\|_F$.
The independent Gaussian entry variances are $t^2/d$ in the first
matrix, $t^2/N$ internally, and $t^2/N^2$ in the readout. Set
$\sigma^2=t^2/N$. The energy is
\begin{equation}\label{eq:energy}
 \kappa_N(W)=\frac dN\|W^1\|_F^2+
       \sum_{s=2}^{L-1}\|W^s\|_F^2+N\|v\|_2^2.
\end{equation}
In energy coordinates the prior is isotropic $N(0,\sigma^2 I)$.
For a consistently labeled sample $S=((x_i,y_i))_{i=1}^s$, condition
only on
\begin{equation}\label{eq:event}
 E_S=\{|z_W(x)|>1/2\ (x\in X_d),\
                         y_i z_W(x_i)>0\ (1\le i\le s)\}.
\end{equation}
Call the posterior $Q_{S,\sigma}$ and define
\[
 e_{S,\sigma}(f)=2^{-d}\bigl|\{x:
       Q_{S,\sigma}(f(x)z_W(x)>0)\le2/3\}\bigr|.
\]
The exponent convention used in the abstract is
$\sigma^2=N^{-\gamma}$, equivalently $t^2=N^{1-\gamma}$.
Thus a fixed $\gamma>1$ is a shrinking-prefactor regime, while
$\gamma=1$ means $t=1$. Weighted depth $L$ includes the readout;
there are $L-1$ hidden layers.

Here $\|g\|_{C^r(\R)}=\max_{0\le j\le r}\sup_{u\in\R}|g^{(j)}(u)|$.

\begin{theorem}\label{thm:main}
Fix $r\ge1$ and $\eta>0$. There exist one fixed computable activation
$\phi$, a computable family $f_d:X_d\to\{-1,1\}$, and constants
$c>0,d_0,D,K$ with the following properties.
\begin{enumerate}[label=(\roman*),leftmargin=*]
\item $\phi$ is odd, real analytic, strictly increasing, and globally
$1$-Lipschitz, with range $(-1,1)$ and $\phi'(0)=1$.
It is strictly concave on $(0,\infty)$ and
$\|\phi-\tanh\|_{C^r(\R)}<\eta$.
\item The targets have teachers of weighted depth
$L_d=Dd+K$, width and coefficient bound $C(d+1)$, and global signed
output margin greater than one. The coefficients can be chosen
with bit lengths polynomial in $d$.
One may take $D=\max(16,r+2)+1$.
\item For every $\zeta>0$ there is $C_\zeta$ such that for
$d\ge d_0$, all $N\ge\exp(C_\zeta d)$, all
$0<\sigma^2\le N^{-(1+\zeta)}$, and every sample of at most
$\exp(cd)$ inputs labeled by $f_d$,
\[
 e_{S,\sigma}(f_d)\ge139/1024>1/8.
\]
\item For each $1\le n\le\exp(cd)$, every fixed finite
$N\ge M(d,n)$ for a polynomial $M$ has the same conclusion
for all sufficiently small $\sigma>0$, uniformly over all samples
of at most $n$ inputs.
\item There is $C_{\rm crit}$ such that the same every-sample
conclusion holds at the actual critical scale $\sigma^2=1/N$
for every $N\ge\exp(C_{\rm crit}d)$ and $d\ge d_0$.
\end{enumerate}
Constants may depend on the fixed $r,\eta$. The targets are
evaluated under the uniform input law.
\end{theorem}
Thus the number of samples is exponential in $L_d$. The finite-scale
width threshold is polynomial in the maximum permitted sample budget,
but is not asserted polynomial in every smaller actual sample size.
The coefficient and architectural teacher resources are polynomial
in depth. The hard information is in the one fixed activation.
The constants may depend on a fixed Gaussian exponent; no uniform
claim for exponents approaching one with the problem resources is
implicit in the theorem.

\paragraph{How the proof fits together.}
The hard tables are selected first, using a circuit-counting argument that
does not refer to the activation. The analytic construction then encodes
those tables, and a small teacher decodes them. The posterior argument
runs in the other direction: sample-consistent memorizers make the
conditioning event sufficiently likely, and a representation ceiling
places most posterior mass on functions unable to approximate the chosen
tables. The cold and fixed-gap scales use a low-product ceiling; the
critical scale uses the separate controlled-path and finite-event
arguments in Appendix~\ref{ac-appendix}. The final conversion from Gibbs
error to confidence failure is the same in both arguments.

\section{Universal capacity and computable hard tables}
The tables must be chosen independently of the activation that later
encodes them. We use a capacity bound allowing different activations
at different neurons.
\begin{proposition}[Universal finite-circuit capacity]\label{prop:capacity}
There are universal constants $A,b$ such that any Boolean function
with a globally unit-margin network of Frobenius product $P$ has a
bounded-fan-in Boolean circuit of size at most
$A(4+\lceil P\rceil)^b$. The network may have arbitrary finite depth
and widths, and each neuron may use its own zero-preserving
$1$-Lipschitz activation. The Boolean depth can be
$O(\log^2(4+P))$; a threshold circuit has depth $O(\log(4+P))$.
\end{proposition}
A proof is given in Appendix~\ref{app:capacity}. Computability of
the activations is unnecessary: the circuits are nonuniform and
store finite tables as constants.

\begin{lemma}\label{lem:hard}
There are universal $a>0,d_0$ and a computable family $f_d$ such
that for $d\ge d_0$, every globally unit-margin network in the class
of Proposition~\ref{prop:capacity} with product at most $e^{ad}$
disagrees with $f_d$ on at least $7/16$ of the uniform cube.
\end{lemma}
\begin{proof}
Enumerate all bounded-fan-in Boolean circuits on $d$ inputs with
at most $\lfloor2^{d/4}\rfloor$ gates. There are
$\exp(O(d2^{d/4}))$ such circuits. A Hamming ball of relative radius
$q=7/16$ among $2^d$-bit truth tables has at most
$2^{H_2(q)2^d}$ elements, with binary entropy $H_2(q)<1$.
The union covers fewer than $2^{2^d}$ tables eventually. Choose
the lexicographically first table at distance at least $q2^d$
from every listed circuit. At smaller dimensions with no such
table choose a constant. Finite exhaustive search computes this
choice. Choose $a>0$ so small that
$A(4+\lceil e^{ad}\rceil)^b\le2^{d/4}$ eventually, and apply
Proposition~\ref{prop:capacity}. The choice refers only to Boolean
circuits, not to the later activation.
\end{proof}


\section{An analytic perturbation with weighted derivative control}
Fix $M=\max(16,r+2)$, $D=M+1$, and $m=M+3$. Let $J$ be a
sufficiently large fixed dyadic integer, increased from $2^{m+10}$
if needed below. Set
\[
 \Delta_d=2^{-d-4},\quad r_d=\Delta_d/J,\quad
 c_{d,k}=-4d-\tfrac14-\Delta_d k,\qquad 0\le k<2^d.
\]
For a sufficiently small fixed positive rational $\tau$, put
$\delta_d=\tau r_d2^{-Md}=\tau(16J)^{-1}2^{-Dd}$ and define
\begin{align}
 S(u)&=\sum_{d\ge1}\sum_{k<2^d}\delta_df_d(k)(u/c_{d,k})^3
                   \exp[-((u-c_{d,k})/r_d)^2],\label{eq:S}\\
 R(u)&=S(u)-S(-u),\qquad \phi(u)=\tanh u+R(u).\label{eq:phi}
\end{align}
\begin{lemma}\label{lem:phi}
The function $S$ is entire and computable. For every fixed
$j\le\max(r,5)$,
\begin{equation}\label{eq:globalder}
 \sup_{\R}|S^{(j)}|\le
 C_j\tau(16J)^{j-1}\sum_{d\ge1}2^{-(M+1-j)d}.
\end{equation}
There are constants $C,C_J$ such that
\begin{align}
 |R'(u)|&\le C\tau\min(u^2,1)\operatorname{sech}^2u,
                                      \label{eq:weightedfirst}\\
 |R''(u)|&\le C_J\tau|\tanh''u|.        \label{eq:weightedsecond}
\end{align}
A sufficiently small $\tau$ makes $\phi$ satisfy all activation
properties in Theorem~\ref{thm:main}.
\end{lemma}
\begin{proof}
On a complex compact set $|\Re z|,|\Im z|\le A$, for all large
$d$ the real part of $(z-c_{d,k})^2$ is at least $c d^2$.
The entire $d$-band is bounded by a polynomial factor times
$2^d\delta_d\exp(-c d^2/r_d^2)$, a summable sequence. This
proves local uniform convergence and entire analyticity.
The explicit tails and computable finite tables also prove
computability, without an efficiency assertion.

For one bump write $v=(u-c)/r$ and $\rho=r/c$.
Its shape is $(1+\rho v)^3e^{-v^2}$ with $|\rho|\le1$.
Its $j$-th derivative is bounded by $C_jr^{-j}e^{-v^2/2}$.
The elementary lattice estimate
\[
 \sup_s\sum_{k\in\mathbb Z}e^{-(s-Jk)^2/2}\le3\qquad(J\ge4)
\]
follows by choosing the nearest lattice point and summing the
remaining tails. This gives \eqref{eq:globalder}.

For a weighted derivative bound use
$\operatorname{sech}^2u\ge e^{-2|u|}$,
$|c_{d,k}|\le4d+5/16$, and
$|u|\le|c_{d,k}|+r_d|v|$.
The derivative's polynomial factor times
$e^{-v^2+2r_d|v|}$ is at most $Ce^{-v^2/2}$, uniformly
since $r_d\le1/128$. Therefore
\[
 \frac{|S'(u)|}{\operatorname{sech}^2u}
 \le C\tau\sum_{d\ge1}e^{-(M\log2-8)d}\le C'\tau.
\]
A second derivative has an extra $r_d^{-1}$, giving
\[
 \frac{|S''(u)|}{\operatorname{sech}^2u}
 \le CJ\tau\sum_{d\ge1}e^{-((M-1)\log2-8)d}\le C_J'\tau.
\]
Both series converge because $M\ge16$.

Near zero the cubic factor gives additional powers.
For $|u|\le1$, a single first derivative divided by $u^2$
is bounded by
\[
 \frac{\delta_d}{|c|^3}
 \left(3+\frac{2|u||u-c|}{r_d^2}\right)e^{-(u-c)^2/r_d^2}.
\]
Here $|u-c|\in[3d,6d]$. Its band sum is at most
$C\tau2^{-(M-1)d}(r_d+r_d^{-1})e^{-9d^2/r_d^2}$.
For a second derivative divided by $u$, the exact bracket is
$\delta_dc^{-3}[6g(v)+6ug'(v)/r_d+u^2g''(v)/r_d^2]$,
where $g(v)=e^{-v^2}$; its band sum is at most
\[
 C\tau2^{-(M-1)d}
   (r_d+r_d^{-1}+r_d^{-3})e^{-9d^2/r_d^2}.
\]
Both series are $O(\tau)$. Since $\operatorname{sech}^2u$
is bounded below on $[-1,1]$, this supplies the additional
$u^2$ in \eqref{eq:weightedfirst}.
Also $|\tanh''u|$ is bounded below by a positive constant
times $|u|$ on this interval, and by a positive constant
times $\operatorname{sech}^2u$ outside it.
Apply these bounds at $u$ and $-u$ to get
\eqref{eq:weightedfirst}--\eqref{eq:weightedsecond}.

The perturbation is odd and vanishes to order three at zero,
so $\phi(0)=0$ and $\phi'(0)=1$.
Choose $\tau$ so both weighted constants are at most
$\varepsilon<1/10$. Then $\phi'(u)>0$, while
\[
 \phi'(u)\le\operatorname{sech}^2u+
       \varepsilon u^2\operatorname{sech}^2u
 \le\operatorname{sech}^2u+\varepsilon\tanh^2u\le1,
\]
using $|u|\le|\sinh u|$.
The second weighted estimate gives $\phi''(u)<0$ for $u>0$.
Uniform summability in \eqref{eq:globalder} implies
$S(u)\to0$ as $|u|\to\infty$, so $\phi$ has limits $\pm1$
and range $(-1,1)$. Shrink $\tau$ once more to obtain the
required $C^r$ distance.
\end{proof}

Two consequences used below, after decreasing $\tau$ if needed, are
\begin{align}
 |\phi(u)-u|&\le2u^2,\label{eq:taylor}\\
 1-\phi(u)&\le Ce^{-2u}\quad(u\ge0).\label{eq:saturation}
\end{align}
The first follows from $\phi'(0)=1$ and a global second-derivative
bound. The second follows by integrating
$\phi'\le(1+\varepsilon)\operatorname{sech}^2$ to infinity.
More precisely $1-\phi(u)$ lies between
$(1-\varepsilon)(1-\tanh u)$ and
$(1+\varepsilon)(1-\tanh u)$.

\section{Finite differences decode the hidden truth table}
Put $h_d=\Delta_d/(2m)$ and define
\begin{equation}\label{eq:decoder}
 F_d(u)=\sum_{j=0}^m(-1)^j\binom mj\phi(u+jh_d).
\end{equation}
\begin{lemma}\label{lem:decode}
For all sufficiently large $d$, uniformly over $0\le k<2^d$,
\[
 F_d(c_{d,k})=\delta_d(f_d(k)+e_{d,k}),\qquad |e_{d,k}|\le1/100.
\]
\end{lemma}
\begin{proof}
The central bump has value exactly $\delta_df_d(k)$.
Every off-center query $u+jh_d$, $1\le j\le m$, lies between
$h_d$ and $\Delta_d/2$ from this center, hence is at least
$h_d$ from all centers in its band.
Gaussian-polynomial absorption and the lattice bound make its
same-band contribution at most
$C\delta_d\exp[-J^2/(16m^2)]$.
Other same-band bumps at the central query contribute at most
$C\delta_de^{-J^2/4}$.
The coefficient absolute sum is $2^m$, so fixed sufficiently
large $J$ makes their aggregate error less than $\delta_d/1000$.

Other-band contributions divided by $\delta_d$ are bounded by
a polynomial factor times
\[
 \sum_{e\ne d}2^{e-D(e-d)}
           \exp[-9(e-d)^2J^2 2^{2e+8}].
\]
This tends to zero. For $e\le d/2$ its exponent dominates
$-c d^2 4^e$; for $d/2<e<d$ it dominates $-c4^{d/2}$;
for $e>d$ it decreases faster still. All prefactors are
dominated in these three splits.
Reflected bumps have distance at least $3(d+e)$ from their
negative-band centers, so satisfy a stronger estimate.

The tanh part is an $m$-th forward difference.
Its iterated-integral formula bounds it by
\[
 h_d^m\sup_{v\in[u,u+mh_d]}|\tanh^{(m)}v|
       \le C_m\Delta_d^m e^{-8d}.
\]
For $v\le-1$, the estimate
$|\tanh^{(m)}v|\le C_me^{2v}$ follows from its convergent
series in $e^{2v}$.
Dividing by $\delta_d=\tau(16J)^{-1}2^{-Dd}$ gives
$O(2^{-2d}e^{-8d})$, since $m=M+3$.
The fixed number of query terms does not affect convergence.
These bounds prove the claim.
\end{proof}

Let $\operatorname{code}(x)=\sum_{i=0}^{d-1}2^i(x_{i+1}+1)/2$.
The first teacher layer has $m+1$ neurons with preactivations
$c_{d,\operatorname{code}(x)}+jh_d$, $0\le j\le m$.
Their biases are $O(d)$ and other coefficients are at most $1/64$.
A second hidden neuron uses the binomial weights in
\eqref{eq:decoder}. Its output has the desired sign and
amplitude $\delta_d(1+O(1/100))$.

Follow it by $Dd+K_0$ scalar hidden layers of weight two.
Choose fixed $K_0$ so $2^{K_0}\tau/(16J)$ is a sufficiently
small positive amplitude $\rho$, up to a factor two.
By \eqref{eq:taylor}, cumulative relative error during the
amplification is $O(\rho)$: the sum of ideal intermediate
amplitudes is a geometric series bounded by $4\rho$.
Choose $\rho$ to make that error less than $1/10$ and use
readout $2/\rho$. The signed output margin exceeds one.
Depth is $L_d=Dd+K_0+3$, active width is $m+1$, and all
coefficients are bounded by $C(d+1)$.
First-layer coefficients are rational with $O(d)$ bits.
Fixed remaining coefficients can be approximated with polynomial
bit lengths while preserving the margin by layerwise Lipschitz
propagation.

\section{A low-energy memorizer in the exponential sample range}
Choose an even integer $A\ge C_1\log(n+2)$ large enough for
\eqref{eq:saturation}, and put
\[
 h_p(x)=\tfrac12\{\phi(A(1-d+p\cdot x))+\phi(A)\}.
\]
The first argument is $A(1-2\operatorname{Ham}(x,p))$, with
absolute value at least $A$. Oddness and saturation imply
$|h_p(x)-\mathbf1_{\{x=p\}}|\le Ce^{-2A}$.
If $p_1,\ldots,p_s$, $s\le n$, are the distinct negative sample
points, then
\[
 z_0(x)=2-4\sum_{j=1}^s h_{p_j}(x)
\]
is within $1/10$ of the ideal memorizer taking values $\pm2$.
The constant feature is $\phi(A)$, with readout
$2/\phi(A)-2s$; each point-feature readout is $-2$.
Thus its first Frobenius norm squared is at most
$F^2=A^2[1+n(d+1)^2]$, and its readout squared norm is
$O((n+1)^2)$.

Zero-pad to
\begin{equation}\label{eq:MB}
 M_0=\lceil\max\{4,n+1,dF^2\}\rceil.
\end{equation}
For $N\ge M_0$, clone every feature $\lfloor N/M_0\rfloor$
times and divide its readout coefficients by that number.
The output is unchanged, the first weighted energy is at
most one, and the old readout $v_0$ has
\begin{equation}\label{eq:B0}
 B=N\|v_0\|_2^2\le B_0
       \le C(d+1)^3(n+1)^3\log^2(n+2).
\end{equation}

Reduce the hard exponent $a$ from Lemma~\ref{lem:hard} so
$a\le1/100$, and fix
\[
 q=\exp(a/(100D))>1.
\]
Choose a fixed small $\alpha>0$
with $\alpha\le c_0(q-1)$ for a sufficiently small universal
$c_0$. Put $\ell=L_d-2$ and
$\rho_1=\alpha q^{1-\ell}$.
The first scalar hidden matrix has every row
$\rho_1v_0^{\mathsf T}$; the remaining $\ell-1$ scalar
hidden matrices have every entry $q/N$.
Their ideal states are
\[
 v_j(x)=\alpha q^{j-\ell}z_0(x).
\]

\begin{lemma}\label{lem:memorizer}
For all sufficiently large $d$ and $n\le\exp(cd)$, where
$c<a/1000$ is fixed, this construction followed by readout
entries $1/(\alpha N)$ has global margin at least $1.5$,
fits every sample, and has energy at most
\begin{equation}\label{eq:K}
 K_d=2+(\ell-1)q^2+\alpha^{-2}\le q^2L_d+C_*.
\end{equation}
The constants are independent of $d,n,N$.
\end{lemma}
\begin{proof}
Equation \eqref{eq:taylor} gives a relative-error bootstrap.
The first actual-to-ideal ratio differs from one by at most
$C|v_1|$. Subsequently write
$y_j=q y_{j-1}(1+\theta_j)$ with
$|\theta_j|\le C|y_{j-1}|$.
Under $|y_i/v_i|\le2$, the sum of these relative errors is at
most $C'\alpha/(q-1)$.
Our choice of $\alpha$ makes this less than $1/100$; the
product of $1+\theta_j$ closes the bootstrap and puts every
ratio in $[9/10,11/10]$.
All preactivations stay in $[-1,1]$.
The final output therefore has margin at least
$(9/10)(1.9)>1.5$, with the signs of $z_0$.

From \eqref{eq:B0}, $\log B_0\le3cd+O(\log d)$, whereas
$q^{-2\ell}=\exp[-ad/50+O(1)]$.
Thus $B\rho_1^2\le1$ eventually, uniformly over the sample
range. This is the first scalar energy.
The other scalar energies are $q^2$ each and the readout
energy is $\alpha^{-2}$; adding the original first energy
proves \eqref{eq:K}.
\end{proof}

The product of the block energies in \eqref{eq:energy} is
$dP_F(W)^2$. Hence
\begin{equation}\label{eq:amgm}
 dP_F(W)^2\le(\kappa_N(W)/L_d)^{L_d}.
\end{equation}
For all large $d$ a globally strict-margin-$1/2$ function with
an energy-$K_d+1$ representative therefore has a unit-margin
representative of product at most
\begin{equation}\label{eq:coldP}
 P_{\rm cold}=\frac2{\sqrt d}
            ((K_d+1)/L_d)^{L_d/2}\le Cq^{L_d}/\sqrt d.
\end{equation}

In the remaining proof write $L=L_d$, $K=K_d$, and $M=M(d,n)=M_0$.
\section{Cold posterior concentration}
\begin{lemma}\label{lem:cold}
For fixed $d,n,L,N$ with $N\ge M(d,n)$,
\[
 \sup_{|S|\le n}Q_{S,\sigma}\{\kappa_N>K+1\}\longrightarrow0
 \qquad(\sigma\downarrow0),
\]
where the supremum is over consistently labeled samples.
\end{lemma}
\begin{proof}
The memorizing witness has margin at least $1.5$ and energy at most $K$.
Continuity of $\phi$ and finiteness of $X_d$ imply that a nonempty
ball in energy coordinates lies inside $E_S$ and has energy at
most $K+1/4$. If $p$ is the parameter dimension, this ball has
Gaussian probability at least
$c_S\sigma^{-p}\exp[-(K+1/4)/(2\sigma^2)]$.
For $\eta=1/[4(K+1)]$, the Gaussian exponential moment gives
\[
 \mu_\sigma\{\kappa_N>K+1\}
 \le\eta^{-p/2}\exp[-(K+3/4)/(2\sigma^2)].
\]
The ratio tends to zero. Only finitely many distinct conditioning
events arise from samples of at most $n$ inputs on the finite
cube, so convergence is uniform over those events.
\end{proof}

\section{Actual finite Gaussian scales}
The following argument uses only $\phi(0)=0$ and
$\Lip(\phi)\le1$. Normalize the input by $\sqrt d$, hidden states
by $\sqrt N$, and use energy-coordinate blocks
\[
 A_1=\sqrt{d/N}\,W^1,\quad A_s=W^s\ (2\le s<L),\quad
 A_L=\sqrt N\,v^{\mathsf T}.
\]
The normalized hidden activation is
$u\mapsto\phi(\sqrt N u)/\sqrt N$, still zero-preserving and
$1$-Lipschitz. The normalized input norm is at most $\sqrt2$.

\begin{lemma}[Uniform perturbation]\label{lem:perturb}
If $\kappa_N(U)\le K+1$ and each normalized block of $B$ has
operator norm at most $\eta$, then
\[
 \sup_{x\in X_d}|z_{U+B}(x)-z_U(x)|
 \le\sqrt2L\eta
 \left(\frac{K+1+2\eta\sqrt{L(K+1)}+L\eta^2}{L-1}
                                                   \right)^{(L-1)/2}.
\]
Define
\begin{equation}\label{eq:eta}
 A_*=\left(\frac{K+2}{L-1}\right)^{(L-1)/2},\quad
 \eta=\min\left\{1,\frac1{4\sqrt{L(K+1)}},
                         \frac1{8\sqrt2LA_*}\right\}.
\end{equation}
Then the output perturbation is at most $1/8$.
\end{lemma}
\begin{proof}
Replace blocks one at a time. Zero preservation bounds intermediate
state norms by the input norm times preceding operator norms.
Lipschitzness bounds each replacement's effect by $\eta$ times
the other $L-1$ factors. Write $a_s=\norm{U_s}_F$; those factors
are at most $a_s+\eta$. The sum of their squares, even including
the omitted factor, is at most
$K+1+2\eta\sqrt{L(K+1)}+L\eta^2$. Arithmetic--geometric mean and
the sum over replacements prove the bound. The second choice in
\eqref{eq:eta} makes this numerator at most $K+2$, and the last
choice gives the error $1/8$.
\end{proof}

Set
\begin{equation}\label{eq:bulk}
 b^2=\max\{2(d+1)\log2+4\log8,\,
                          16\log9+8\log(8L)\}.
\end{equation}
In energy coordinates let $T$ be the set where the first and last
blocks have Frobenius norm at most $tb$, and every internal block
has operator norm at most $tb$. This set is closed, convex, and
centrally symmetric. It has raw Gaussian mass at least $1/2$.
Indeed the first and last bounds each fail with probability at
most $1/8$ by the chi-square exponential moment. Two $1/4$-nets
of the $N$-dimensional unit sphere have at most $9^{2N}$ pairs;
the scalar Gaussian tail and a union bound give, for each internal
block, failure probability at most
\[
 2\exp(2N\log9-Nb^2/8)\le1/(4L).
\]
The sum of all these failure bounds is at most $1/2$.

\begin{proposition}\label{prop:finite}
Suppose $tb\le\eta$ for \eqref{eq:eta}--\eqref{eq:bulk}.
For every sample of at most $n$ inputs, the posterior gives mass
at least $1-2e^{-1/(2\sigma^2)}$ to Boolean functions admitting a
global unit-margin representative of product at most
\begin{equation}\label{eq:Pfin}
 P_{\rm fin}=\frac8{3\sqrt d}((K+1)/L)^{L/2}\le Cq^L/\sqrt d.
\end{equation}
\end{proposition}
\begin{proof}
Let $w_S$ be the memorizing witness.
Lemma~\ref{lem:perturb} gives $w_S+T\subset E_S$.
The Gaussian shift formula and central symmetry of $T$ imply
\[
 \mu_\sigma(E_S)\ge\mu_\sigma(w_S+T)
       \ge\tfrac12\exp[-K/(2\sigma^2)].
\]
Gaussian isoperimetry for a set of Gaussian mass at least $1/2$
gives
\[
 \mu_\sigma\{\dist(W,T)^2>K+1\}
        \le\exp[-(K+1)/(2\sigma^2)].
\]
Distance here is Euclidean distance in energy coordinates.
Division bounds the posterior exceptional mass by
$2e^{-1/(2\sigma^2)}$.

Outside this set let $B$ be the metric projection of $W$ onto
$T$ and set $U=W-B$. Then $\kappa_N(U)\le K+1$, and all
normalized blocks of $B$ have operator norm at most $tb\le\eta$.
Lemma~\ref{lem:perturb} makes $z_U$ differ from $z_W$ by at most
$1/8$ everywhere. Their signs agree and $|z_U|>3/8$.
Multiply its readout by $8/3$, and apply \eqref{eq:amgm} and
\eqref{eq:K}, to obtain \eqref{eq:Pfin}.
\end{proof}


Fix $\zeta>0$. It suffices to require
\begin{equation}\label{eq:width}
 N\ge\left\lceil\max\{M,2,(b/\eta)^{2/\zeta},2\log64\}\right\rceil.
\end{equation}
Then $\sigma^2\le N^{-(1+\zeta)}$ implies
$t=\sqrt{N\sigma^2}\le N^{-\zeta/2}$, so $tb\le\eta$;
the posterior exceptional mass is at most $1/32$.
For $K\le q^2L+C_*$,
\[
 ((K+2)/L)^L\le Cq^{2L},\qquad
 A_*^2\le e((K+2)/L)^L\le Cq^{2L}.
\]
Thus $\eta^{-1}\le C\poly(L)q^L$.
The bulk constant is polynomial in $d,L$, and
$M\le\exp(Cd)$ uniformly for $n\le\exp(cd)$.
Since $L=Dd+O(1)$, a bound $N\ge\exp(C_\zeta d)$
satisfies \eqref{eq:width} throughout this sample range.

\section{Completion at the cold and fixed-gap scales}
Since $L_d\log q=ad/100+O(1)$, both
\eqref{eq:coldP} and \eqref{eq:Pfin} are at most $e^{ad}$
for all sufficiently large $d$.
Lemma~\ref{lem:hard} gives error at least $7/16$ for every
Boolean function in those posterior events.
By the cold comparison or finite-scale projection argument,
the event has posterior mass at least $31/32$.
The posterior Gibbs error is at least
$(31/32)(7/16)=217/512$.
If $e=e_{S,\sigma}(f_d)$, the Gibbs error is at most
$(1-e)/3+e$. Hence
\[
 e\ge[3(217/512)-1]/2=139/1024>1/8.
\]
All estimates are uniform over sample realizations. Together
with the teacher construction this proves the cold and fixed-gap
parts of Theorem~\ref{thm:main}.

\section{The critical endpoint}
The same activation satisfies the hypotheses of the critical
target-mass and canonical-state results proved in
Appendix~\ref{ac-appendix}. Here we check the activation-specific
strict potential explicitly.
For $\tau$ sufficiently small, \eqref{eq:weightedfirst} gives
\[
 |R(u)|\le\varepsilon\min(|u|^3/3,1).
\]
Both $\phi$ and tanh are bounded by one and by $|u|$, so
$|\phi(u)^2-\tanh(u)^2|\le2\varepsilon\min(u^4,1)$.
For standard Gaussian $Z$ and $0\le q\le2$,
\[
 \left|\E e^{\phi(\sqrt qZ)^2}
             -\E e^{\tanh(\sqrt qZ)^2}\right|
       \le6e\varepsilon q^2.
\]
There is $c_0>0$ with
$\log\E e^{\tanh(\sqrt qZ)^2}\le q-c_0q^2$ on $[0,2]$:
for $G(u)=e^{\tanh^2u}$, direct differentiation gives
$G''<2G$ away from zero, and the heat equation gives strict
inequality at positive $q$; Taylor expansion at zero gives
$q-q^2+O(q^3)$. Compactness supplies $c_0$.
Taking $\varepsilon\le c_0/(12e)$ proves
\begin{equation}\label{eq:strictpotential}
 \log\E e^{\phi(\sqrt qZ)^2}\le q-(c_0/2)q^2
                         \qquad(0\le q\le2).
\end{equation}
We may impose this further smallness on $\tau$ without changing
any preceding argument.
The global derivatives through order five are bounded by
\eqref{eq:globalder}, and the weighted second-derivative estimate
at zero gives $\phi'''(0)=-2+O(\varepsilon)<0$.
That estimate
also gives $|\phi''(u)|\le3\phi(u)$ for $u>0$, after decreasing
$\tau$: integrate $|R'|\le\varepsilon\operatorname{sech}^2$
to get $\phi(u)\ge(1-\varepsilon)\tanh u$ and use
$|\phi''|\le(1+\varepsilon)|\tanh''|$.
All hypotheses in Appendix~\ref{ac-appendix} are therefore met.

For the sample memorizer
$g_S=1-2\sum_{p\ {\rm negative\ in}\ S}\mathbf1_{\{p\}}$,
Fourier variation is at most $1+2n$.
The appendix proves that whenever $L\ge C\log(d+n+L)$,
at the actual scale $\sigma^2=1/N$ and every
$N\ge2^{Cd}(d+n+L)^C$, the posterior gives mass at least
$31/32$ to functions with Boolean circuits of size
$\poly(d+n+L)$, uniformly over every sample of at most $n$
inputs.
Reduce $c>0$ in Theorem~\ref{thm:main} so that
$n\le e^{cd}$ and $L_d=Dd+K$ imply both the depth condition
and this circuit-size bound being at most $2^{d/4}$.
The definition of $f_d$ then gives error at least $7/16$
for each such function. The same Gibbs calculation yields
$e_{S,\sigma}(f_d)\ge139/1024$.
Finally $2^{Cd}(d+n+L_d)^C\le e^{C'd}$ in this range.
This proves the critical part and completes the theorem.

\section{Why the activation cannot be efficiently evaluated}
\begin{proposition}\label{prop:evaluation}
There are fixed constants $C,c'>0$ such that approximating this
$\phi$ to error $2^{-Cd}$ on all dyadic inputs of magnitude
at most $Cd$ and bit length at most $Cd$ cannot be done by
Boolean circuits of size less than $\exp(c'd)$ for all
sufficiently large $d$. In particular no algorithm performs
these evaluations in time polynomial in the input length
and requested precision.
\end{proposition}
\begin{proof}
Round the teacher's rational parameters to dyadics with $O(d)$
bits. Its one binomial-combination layer has a fixed norm,
and its remaining amplifying layers have scalar weight two.
Layerwise Lipschitz propagation therefore bounds the output
error from parameter, arithmetic, and activation-evaluation
errors by $C_1d\,2^{Dd}\delta$ when every primitive error is
at most $\delta$. Taking $\delta=2^{-C_2d}$ with fixed
sufficiently large $C_2$ preserves the global output signs.
All preactivation arguments have magnitude $O(d)$ and
dyadic bit length $O(d)$. There are polynomially many
arithmetic operations and activation calls.

If one accurate activation evaluator had circuit size
$\exp(c'd)$, replacing all activation calls by copies of it
would give a circuit for $f_d$ of size at most
$\poly(d)\exp(c'd)$. Choosing fixed $c'<(\log2)/4$
makes this smaller than $2^{d/4}$ eventually, contrary to
the defining circuit hardness of $f_d$.
Enlarge $C$ to cover the required argument and precision
constants. Polynomial running time would in particular
give polynomial-size circuits, so is excluded as well.
\end{proof}

\section{Scope}

The activation is fixed for the entire dimension sequence and
is computable by finite search and convergent analytic series.
It is intentionally computationally contrived. The proof
asserts no efficient evaluator for it: inverse-exponential
accuracy at its lookup bands reveals circuit-hard truth tables.
The teacher architecture, coefficient sizes, and coefficient
precision remain polynomial in its depth.

There is no limiting inference to the exact tanh activation.
Decreasing the perturbation size decreases the encoded signal
and increases the fixed number of amplification layers.
The decoder deliberately amplifies that small activation difference
to a constant output margin; its constants are not uniform as the
perturbation tends to zero.

Thus oddness, analyticity, monotonicity, concavity, exact slope
one at zero, and arbitrary closeness to tanh in any fixed
finite number of derivatives do not by themselves give a
universal polynomial-depth learning theorem.
The exact tanh activation is unchanged by this statement.
No lower bound for the exact tanh, ReLU, or GELU activation
follows from this construction.

\appendix
\section{Universal finite-circuit capacity}\label{app:capacity}
We prove Proposition~\ref{prop:capacity}, including its parallel-depth
assertion. All activations in this appendix may differ by neuron.

\paragraph{Relevant coordinates.}
Let $T$ be the Frobenius product after the first matrix. If coordinate
$i$ is relevant, a bichromatic cube edge has output difference at
least two, so $T\norm{W^1e_i}\ge1$. Summing the squares shows that
the number $k$ of relevant coordinates is at most $P^2$.
Freeze the others to signs without increasing the first Frobenius
norm: the expected squared norm of the new bias column under
independent random signs is the sum of the old bias and frozen-column
squared norms. Some assignment attains at most this expectation.
Put $R=\sqrt{k+1}$ and $B=4+\lceil P\rceil$. Also $P\ge1$, since
zero preservation and the margin imply
\[
 1\le\E z(x)^2\le T^2\E\norm{W^1(1,x)}^2
          =T^2\norm{W^1}_F^2\le P^2.
\]
In particular $R\le B$.

\paragraph{A short prefix can be compressed.}
For adjacent matrices $W,V$, let $w_i$ be row $i$ of $W$ and
$v_i$ column $i$ of $V$, with $a_i=\norm{w_i}$, $b_i=\norm{v_i}$,
and $A_0=\sum_i a_ib_i\le\norm W_F\norm V_F$.
Delete zero contributions. Normalize the incoming rows to
$u_i=w_i/a_i$ and activations to
$\psi_i(s)=\phi_i(a_i s)/a_i$.
Sample $m$ indices with probabilities $p_i=a_ib_i/A_0$.
Each sampled neuron has row $u_i$, activation $\psi_i$, and
outgoing column $A_0v_i/(mb_i)$. These activations remain
zero-preserving and $1$-Lipschitz. For every draw the new adjacent
Frobenius product is exactly $A_0$, and its expected contribution
to the next preactivation is the old one.

Let $S$ be the preceding product. The incoming state norm is at
most $SR$, and an unaveraged sampled vector has norm at most
$A_0SR$. Hilbert-space bounded-difference concentration and a
union bound over $X_k$ give some draw with uniform preactivation
error $O(A_0SR\sqrt{\log(8|X_k|)/m})$.
The suffix after $V$ has Lipschitz constant at most
$P/(S\norm W_F\norm V_F)$. Thus final output error is at most
$O(PR\sqrt{\log(8|X_k|)/m})$.
For $h=O(\log B)$ choose
$m=C(PRh)^2\log(8|X_k|)$, and perform this replacement through
the first $\min\{h,L-1\}$ hidden layers. The product never
increases, and sufficiently large universal $C$ gives total
output error at most $1/4$. Double the readout: the resulting
network has margin at least $3/2$, product $P_0\le2P$, and
retained widths polynomial in $B$. A following unretained matrix
may still have arbitrarily many rows, which does not affect these
claims.

\paragraph{Two-ball refinement.}
The image of a radius-$r$ ball under a matrix $W$ is covered by at
most two balls of radius $r\norm W_F/\sqrt2$.
Normalize $\norm W_F=1$ and let $a$ be its leading singular value.
If $a^2\le1/2$, one ball at zero works. Otherwise put
$b^2=1-a^2$, $c=(a^2-b^2)/(2a)$, and use centers $\pm cu$ along
the leading left-singular direction. For a unit input with
leading right-singular coordinate $t$, the squared distance to
the matching center is at most
\[
 (a|t|-c)^2+b^2(1-t^2).
\]
It is convex in $|t|\in[0,1]$, and both endpoint values are
$1/(4a^2)\le1/2$. Scaling and translation prove the claim.
Coordinatewise Lipschitz activations preserve the covering radii.

Let $S_j$ be the product of the first $j$ modified matrix norms.
Iteration gives at most $2^j$ balls of radius $RS_j2^{-j/2}$.
Replace each nonempty ball center by an actual hidden image,
doubling its radius. Choose
$h=\lceil2\log_2(64B^2)\rceil$.
If at least $h$ hidden layers exist, normalize their final state
by $S_h$. Every input is within $1/(8\overline P)$ of an actual
representative, where $\overline P=4B$. Opposite-label
representatives are at distance at least $2/\overline P$:
the normalized suffix is $P_0$-Lipschitz and both outputs have
margin at least one. There are at most $2^h=\poly(B)$ centers.
If fewer hidden layers exist, retain the entire compressed
network and its readout instead.

\paragraph{Normalize before rounding.}
Put $M_s=\norm{W^s}_F$, $S_s=\prod_{i\le s}M_i$,
$A_s=W^s/M_s$, and
$\Psi_{s,i}(u)=\phi_{s,i}(S_su)/S_s$. No $M_s$ vanishes,
since that would force all outputs to vanish. The normalized
states satisfy
\[
 u_s=\Psi_s(A_su_{s-1}),\qquad \norm{A_s}_F=1,\qquad
 \norm{u_s}\le R.
\]
Every $\Psi_{s,i}$ is zero-preserving and $1$-Lipschitz.
Choose inverse-polynomial meshes in $B$ for the matrix entries
and for scalar tables of $\Psi_{s,i}$ on $[-R-2,R+2]$.
If matrix Frobenius error is at most $\eta$ and scalar lookup
error at most $2\delta$, the state error obeys
\[
 e_s\le(1+\eta)e_{s-1}+\eta R+2\sqrt C\,\delta,
\]
where $C=\poly(B)$ bounds the retained widths.
Taking $\eta,\delta$ below suitable inverse polynomials in
$B,C,h$ gives $e_s<1/(256\overline P)$ throughout.
Approximate affine inputs remain in the stored interval.
Store the normalized representatives to the same Euclidean
accuracy. Squared-distance comparisons at threshold
$1/\overline P^2$ distinguish a close correct-label center from
every opposite-label center. The output is positive precisely
when a positive-labeled center lies below that threshold.
In the short-network case, round the normalized readout as
well; its margin at least $1/\overline P$ dominates these errors.

All tables have polynomial length in $B$, with $O(\log B)$ bits
per entry. Fixed-precision sums, products, comparisons, and
lookups therefore give a Boolean circuit of size $\poly(B)$
with universal degree. Carry-save addition trees followed by
carry propagation take $O(\log B)$ depth per retained layer,
giving total $O(\log^2 B)$ Boolean depth.
Alternatively each fixed affine comparison of preceding state
bits is a threshold gate. Interval tests and table-bit selection
take constant threshold depth per layer. Squaring the final
$O(\log B)$-bit differences uses polynomial-size lookup tables,
and threshold gates sum the squares. The threshold depth is
$O(\log B)$. This proves Proposition~\ref{prop:capacity}.

\section{Critical target mass and posterior circuits}
\label{ac-appendix}

This appendix proves the finite-width critical input used in the main
theorem. Throughout, constants may depend on one fixed activation $\phi$.
We assume that $\phi$ is odd, strictly increasing, bounded by one,
and $1$-Lipschitz; that it is $C^5$ with bounded derivatives through
order five, concave on $(0,\infty)$, and satisfies
\[
 \phi'(0)=1,\qquad \phi'''(0)<0,\qquad
 \phi'(u)\ge c_0\operatorname{sech}^2u.
\]
Writing $\beta=\lim_{u\to\infty}\phi(u)>0$, assume
$\beta-\phi(u)\le C(1+u)^{-a}$ for $u\ge0$ and some $a>0$.
The final assumption is the strict potential
\begin{equation}\label{ac-potential}
 \log\E e^{\phi(\sqrt q Z)^2}\le q-cq^2,
                  \qquad 0\le q\le2,
\end{equation}
where $Z$ is standard Gaussian. The main text verifies all these
conditions for the activation constructed there.

Set $H=L-1\ge1$. For a finite input set $X$, a controlled critical
path starts at $Q_0(x,y)=(1+x\cdot y)/d$ and consists of
\[
 \rho_j\ll\mathcal N(0,Q_{j-1}),\quad
 D_j=\operatorname{KL}(\rho_j\Vert\mathcal N(0,Q_{j-1})),\quad
 Q_j(x,y)=\E_{\rho_j}\phi(G_x)\phi(G_y).
\]
Singular references are permitted. Its output $z\in\operatorname{range}Q_H$
has readout cost $r=z^{\mathsf T}Q_H^\dagger z/2$, and total action
$A=\sum_jD_j+r$. Let $J_L(f)$ be the infimum of this action over
full-cube paths with $f(x)z(x)\ge1/2$ everywhere. These definitions
will be related to actual finite-width probabilities below.

\subsection{Scalar control, cooling, and a target program}

For $q>0$ let $V(q)$ be the infimum, over finite scalar paths starting
at $q_0=q$, of
\[
 \sum_{j=1}^m\operatorname{KL}(\eta_j\Vert\mathcal N(0,q_{j-1}))
       +\frac1{8q_m},\qquad q_j=\E_{\eta_j}\phi(G)^2.
\]
The length-zero path is included.

\begin{lemma}[Scalar baseline and short paths]\label{ac-scalar}
The function $V$ is nonincreasing and has a finite positive limit
$V_*:=\lim_{q\downarrow0}V(q)$. For a fixed $C$,
\[
 0\le V_*-V(q)\le Cq\quad(0<q\le2).
\]
A scalar transition of entropy $D$ from variance $q$ to $q'$ satisfies
\[
 D\ge V(q)-V(q'),\qquad D\ge q'-q+cq^2\quad(q\le2).
\]
For every sufficiently small $p>0$ there is a strictly increasing
scalar path starting at $p$, of total cost at most $V(p)+p$, whose
length is at most $C/p$ and for which
\begin{equation}\label{ac-short-path}
 \sum_j q_j\le C\log(2/p),\qquad
 \sum_j\sqrt{D_j}\le C\log(2/p).
\end{equation}
\end{lemma}
\begin{proof}
Rescaling a controlled preactivation when its reference variance
increases preserves its Gaussian relative entropy and increases its
$\phi$-square moment. Applying this successively proves monotonicity.
Oddness, $\phi'(0)=1$, and bounded third derivatives give
$F(u):=\E\phi(\sqrt uZ)^2\ge u-Cu^2$ for small $u$; outside a
fixed neighborhood of zero use Gaussian tails and $|\phi(z)|\le|z|$.
Inflating the reference variance by $1+Kq$, with fixed large $K$,
therefore increases a small $q$ by between $c_1q^2$ and $C_1q^2$
at Gaussian cost at most $C_2q^2$. Continuous adjustment of the
inflation hits any prescribed small endpoint $p$ exactly. Telescoping
the cost-to-increment ratio gives a bridge from arbitrarily small
positive variance to $p$ costing at most $Cp$. A direct readout at
one fixed $p$ proves finiteness of $V_*$, and appending a nearly
optimal suffix proves $V_*-V(p)\le Cp$. Boundedness extends this
inequality to $q\le2$. For starting $q\le1$, every final variance,
including that of a length-zero path, is at most one; hence $V_*\ge1/8$.
The Bellman inequality follows by appending a nearly optimal suffix,
and entropy duality applied to~\eqref{ac-potential} gives dissipation.

Delete each nonincreasing transition of a near-optimal path and rescale
the subsequent controls as above. Cost does not increase, so the path
may be taken strictly increasing. Its cost is bounded by a fixed
constant $B$. Gaussian entropy duality and $\phi(z)^2\le z^2$ bound
each ratio $q'/q$ by a constant depending on $B$. The entropy spent
before a state $q$ is at most $p+Cq$, by the scalar suffix bound.
In a dyadic band $u\le q<2u$, its expenditure, including its exit
transition, is therefore $O(u)$: the exit state is at most a fixed
multiple of $u$. Summed dissipation bounds its length by $O(1/u)$,
apart from at most one exit transition. The band contributes $O(1)$
to both sums in~\eqref{ac-short-path}, using Cauchy--Schwarz for the
second. Summing the geometric length bounds and logarithmically
many bands proves the assertion. States above a fixed small constant
occupy only a bounded number of transitions by the same dissipation.
\end{proof}

\begin{lemma}[An exact low-cost Fourier seed]\label{ac-seed}
Let $A_F(f)=\sum_T|\widehat f(T)|$. There is a first-row controlled
law, with auxiliaries and a score $b_1$, such that
\[
 \E b_1^2=1,\quad \E[b_1\phi(w\cdot(1,x))]=\sqrt{v_1}f(x),
 \quad \E\phi(w\cdot(1,x))^2=q_1
\]
with $q_1$ independent of $x$, relative entropy at most $1/H$, and
$v_1^{-1},\norm{b_1}_\infty$ bounded by a fixed polynomial in
$(d+H+2)A_F(f)$. The reference weight marginal is exactly
$\mathcal N(0,I/d)$.
\end{lemma}
\begin{proof}
For centered Gaussian $G$ of variance at most two, split at $G=-M/2$
to obtain
\[
 \beta-\E\phi(M+G)\le CM^{-a}+Ce^{-M^2/16}.
\]
A parity on $k$ coordinates is a signed sum of $k$ threshold steps
in their Hamming sum and a constant, with coefficient magnitudes
at most one. Each step has an affine argument separated from zero
by one. Shift first-row Gaussian means by a polynomially large
multiple of these affine coefficients, and mix the signed detectors;
reflection of the row implements a negative sign. Their normalized
mixture has mean $g$ with $\chi_T(x)g(x)\ge c/(k+1)$ everywhere.
Every component costs at most $D_* =\poly(d)$ in Gaussian KL.

Average this law over cube translations $x\mapsto s x$, reflecting
the output row by $\chi_T(s)$. Its mean becomes exactly
$c_T\chi_T(x)$, where $c_T=\E_x\chi_T(x)g(x)\ge c/(d+1)$.
Its second-moment diagonal is constant. The transformations preserve
the Gaussian reference, and entropy convexity preserves the cost
$D_*$. Mix these laws, with reflections according to
$\operatorname{sign}\widehat f(T)$, in proportions
$|\widehat f(T)|/c_T$. The resulting law has mean $f/Z_f$,
constant diagonal, cost at most $D_*$, and
$Z_f\le C(d+1)A_F(f)$.

Dilute this law into the free first-row law with probability
$\theta=(HD_*)^{-1}$, increasing $D_*\ge1$ if needed. Retain its
Bernoulli tag $B_0$ and use $b_1=B_0/\sqrt\theta$. The coefficient
is $\sqrt\theta f/Z_f$, so $v_1=\theta/Z_f^2$. In the reference
use the same tags and free Gaussian weights in every branch. Its
weight marginal is the original Gaussian, and the joint entropy is
at most $\theta D_*\le1/H$. All claimed bounds follow. In particular,
no entropy proportional to the number of translation or Fourier tags
has been charged.
\end{proof}

\begin{lemma}[A scalar-channel target program]\label{ac-program}
Put $R=(d+L+2)A_F(f)$ and $S=1+\log R$. If $H\ge CS$, there is
an exact $H$-layer controlled program with scores $b_j$ and deterministic
positive components $v_j$, residual diagonals $r_j\ge0$, such that
\[
 \E b_j^2=1,\quad \E[b_jh_j(x)]=\sqrt{v_j}f(x),\quad
 \E h_j(x)^2=v_j+r_j\le1.
\]
Its charged scalar costs $D_j$ obey
\begin{equation}\label{ac-program-bounds}
 A:=\sum_jD_j+\frac1{8v_H}\le V_*+C S^2/H,
 \quad K:=\sum_{j=2}^H\sqrt{D_j}\le CS,
 \quad |b_j|\le B\le e^{CS}.
\end{equation}
For $j\ge2$, the controlled row uses a scalar $Z$, of entropy cost
at most $D_j$ relative to a reference with standard Gaussian $Z$
marginal, and independent Gaussian residual noise. At one input its
feature is
$\phi(\sqrt{v_{j-1}}f(x)Z+\sqrt{r_{j-1}}G)$.
The score depends on $Z$ and auxiliaries, not on $G$.
\end{lemma}
\begin{proof}
We supply the scalar estimates and the depth accounting. Write
$b(u)=\E\phi'(\sqrt uZ)$. Taylor expansion, with bounded fifth
derivatives and Gaussian tails, gives
\begin{equation}\label{ac-cooling-expansions}
 F(u)=u+\phi'''(0)u^2+O(u^3),\quad
 b(u)=1+\tfrac12\phi'''(0)u+O(u^2),\quad
 \frac{u b(u)^2}{F(u)}=1+O(u^2).
\end{equation}
Also $F(u)\le u-cu^2$ by~\eqref{ac-potential}. In a free layer,
Gaussian integration by parts gives at least $v b(q)^2$ as the
component generated from $vff^{\mathsf T}$ at diagonal $q$.
The component-to-diagonal ratio therefore loses factors
$q b(q)^2/F(q)$. Their logarithms are $O(q^2)$ near zero, and
they have a positive lower bound on each remaining compact interval.
Since free dissipation bounds $\sum q_j^2$, cooling to the first
diagonal at most $\delta$ takes $O(1/\delta)$ layers and retains
component at least $c\delta v_1$. If already below $\delta$, skip
cooling; the same lower bound holds. Cooling costs zero entropy.

The estimate
\begin{equation}\label{ac-derivative-comparison}
 1-\phi'(z)\le C\phi(z)^2
\end{equation}
follows near zero by odd $C^3$ regularity and away from zero by
positivity of $\phi$. Inflate only the target Gaussian coordinate
by variance factor $1+e$. Its cost is
$[e-\log(1+e)]/2\le e^2/4$. If the current component and residual
diagonal are $v,r$ and $q=v+r$, the full conditional-mean component
and new residual satisfy
\begin{equation}\label{ac-growth}
 v'\ge(1+e)v b(q+ev)^2
       \ge(1+e)v[1-C(q+ev)],\qquad r'\le r
\end{equation}
in the small-variance range. The first inequality is the first Hermite
projection; the last follows from conditional Gaussian Poincare.
The component is continuous and increasing in $e$, and at $e=0$
it is at most $v$. Thus at $q\le\delta$ one can hold $v$ exactly
with $e=O(\delta)$, at cost $O(\delta^2)$, preserving $q\le\delta$.
If $r\le\delta$ and $v\le p$, one can prescribe
\[
 v' =\min\{p,v[1+c_1(\delta+v)]\}
\]
with $e=O(\delta+v)$ and cost $O((\delta+v)^2)$. From
$v_0=\delta e^{-C_1S}$ this takes $O(S/\delta)$ steps and costs
$O(p+S\delta)$. Indeed, on full steps,
$(\delta+v)^2=O(\Delta v+\delta\Delta v/v)$;
their square-root costs sum to $O(\log(p/v_0)+1)=O(S)$ when
$p/\delta$ is fixed. The prescribed sequence fixes the length
independently of the residual and any preceding holding time.

At $v=p$, $r\le\delta\le p$, hold the full component exactly.
The total Gaussian preactivation variance lies in $[p,Cp]$, and
\[
 \E\phi'(\sqrt uZ)^2=1+\phi'''(0)u+O(u^2)\le1-cp.
\]
The derivative form of conditional Poincare now gives
$r'\le(1-cp)r$. Thus $O(\log(H+2)/p)$ purification stages leave
$r\le\delta H^{-4}$, at cost $O(p\log(H+2))$ and square-root
cost $O(\log(H+2))$.

For the scalar suffix, $|\phi''(z)|\le C\phi(z)$ for $z>0$:
use oddness and $C^3$ regularity near zero, positivity on compact
intervals, and bounded $\phi''$ elsewhere. If $P_r$ denotes Gaussian
heat convolution, then
\begin{equation}\label{ac-heat}
 e^{-Cr/2}\phi(z)\le P_r\phi(z)\le\phi(z),\qquad z\ge0.
\end{equation}
Indeed $g=\phi''+C\phi$ is odd and nonnegative on the positive
half-line. Pairing the Gaussian density at $y,-y$ shows
$P_sg(z)\ge0$ for $z\ge0$, so
$\partial_s[e^{Cs/2}P_s\phi(z)]\ge0$. The upper bound follows
in the same way from concavity and the odd function $\phi''$.
Concavity also gives $\phi(\alpha z)\ge\alpha\phi(z)$ for
$0\le\alpha\le1$ and $z\ge0$.

Lift the short scalar path of Lemma~\ref{ac-scalar} starting at $p$.
For a nominal law $\eta$ from variance $q$ and an actual component
$v\le q$, use physical scalar $\sqrt{v/q}X$, $X\sim\eta$.
Its Gaussian relative entropy is unchanged. Equation~\eqref{ac-heat}
and concavity imply
\[
 e^{-Cr}(v/q)\le v'/q'\le1,
\]
and the residual remains at most $r$. Starting at $v=q=p$, the
lift of $m\le C/p$ stages therefore retains at least
$e^{-Cm\delta H^{-4}}$ of every nominal component. Since
$\delta\le p$, all components stay at least $p/2$, and the extra
readout cost is $O(H^{-4})$. Its charged square-root cost is
$O(\log(2/p))$ by~\eqref{ac-short-path}.

Choose $\delta=B_1S/H$ and $p=B_2S/H$, with fixed sufficiently
large $B_2\ge B_1$, and then increase the constant in $H\ge CS$
so both scales are small. After cooling reduce the component to
$v_0=\delta e^{-C_1S}$, which is permitted by Lemma~\ref{ac-seed}.
Cooling, prescribed growth, purification, and the chosen scalar
suffix occupy at most $H-1$ stages after choosing $B_1,B_2$.
Insert exactly the unused number of holding stages before growth.
The total cost is
\[
 A\le H^{-1}+V(p)+p+
 C[H\delta^2+p+S\delta+p\log(H+2)+H^{-4}]
 \le V_*+CS^2/H.
\]
Holding contributes $O(H\delta)=O(S)$ to the square-root sum;
the other contributions were bounded above, proving $K\le CS$.
This is a bound on the charged costs themselves, not an inference
from smaller marginal entropies.

The full conditional-mean score is
$P_r\phi(\sqrt vZ)/\sqrt{v'}$, of unit second moment. To retain
a smaller prescribed component $0<v''\le v'$, replace it by
\[
 \sqrt{v''/v'}\,b+\sqrt{1-v''/v'}\,U,
\]
where $U$ is an independent random sign with the same reference law.
This has unit second moment, the desired coefficient, and absolute
bound $1+(v'')^{-1/2}$, with no extra entropy. Every retained component
in the construction is at least $e^{-CS}$, proving the score bound.
At each stage an independent Gaussian residual vector with the current
residual covariance realizes the stated full-cube law. Only its
constant diagonal enters the one-input identities.
\end{proof}

\subsection{Finite-width realization without a rate-limit argument}

\begin{lemma}[Finite scalar-channel probability]\label{ac-channel}
Suppose a program as in Lemma~\ref{ac-program} has action $A$,
distortion $K=\sum_{j\ge2}\sqrt{D_j}$, and score bound $B\ge1$.
For every $\tau>0$ and every integer
\begin{equation}\label{ac-channel-width}
 N\ge CB^4H^4e^{4K}(A+1)^6(1+\tau^{-1})^4
       [\log(2|X|)+\log(H+A+\tau^{-1}+3)],
\end{equation}
the original critical Gaussian prior satisfies
\[
 \mu_N\{f(x)z_W(x)>1/2\ (x\in X)\}\ge e^{-N(A+\tau)}.
\]
The lemma uses only boundedness and the unit Lipschitz constant of the
activation, in addition to the specified program.
\end{lemma}
\begin{proof}
For $N$ independent score/feature rows set
\[
 M=N^{-1}\sum_i b_i^2,\quad
 \widehat c(x)=\frac{N^{-1}\sum_i b_ih_i(x)}{\sqrt M},\quad
 \widehat q(x)=N^{-1}\sum_i h_i(x)^2,\quad
 \widehat s(x)=\sqrt{\widehat q(x)-\widehat c(x)^2}.
\]
If the three raw moments have errors at most $\alpha\le1/4$,
then, writing $c=\E bh$ and $s=\sqrt{\E h^2-c^2}$,
\begin{equation}\label{ac-score-metric}
 \max_x\sqrt{(\widehat c-c)^2+(\widehat s-s)^2}
                         \le8\sqrt\alpha.
\end{equation}
For example $|\widehat c-c|\le4\alpha$ and the squared residual
error before taking square roots is at most $13\alpha$.
The coefficients $a_i=b_i/(N\sqrt M)$ have norm squared $1/N$;
the empirical kernel decomposes exactly as
$K_{\rm emp}=\widehat c\widehat c^{\mathsf T}+R_{\rm emp}$ with
$R_{\rm emp}\succeq0$.

Condition on the entire preceding construction and set $u=\sqrt N a$.
For a fresh row draw its prescribed scalar $(Z,U)$ and independent
standard Gaussian $\xi\in\R^N$, and use
\begin{equation}\label{ac-row-realization}
 w=[Zu+(I-uu^{\mathsf T})\xi]/\sqrt N.
\end{equation}
Under the reference, $Z$ is standard Gaussian, so the incoming
weight marginal is exactly $\mathcal N(0,I/N)$. Its controlled
relative entropy is at most $D_j$ by data processing. At one input
the preactivation has law $\widehat c(x)Z+\widehat s(x)G$,
with independent standard $G$. Couple this to the ideal preactivation
using the same $Z,U,G$. Gaussian entropy duality gives
$\E Z^2\le e^{2\sqrt{D_j}}$, so the preactivation and feature RMS
errors are at most $e^{\sqrt{D_j}}$ times the preceding metric error.
Orthogonal projection onto the common unit score, followed by the
reverse triangle inequality for residual norms, gives the same bound
for the population coefficient/residual metric. Fresh sampling and
\eqref{ac-score-metric} consequently give
\begin{equation}\label{ac-metric-recursion}
 e_H\le8He^K\sqrt\alpha.
\end{equation}
Only one-input comparison couplings are used; the sampled network
keeps its usual joint residual Gaussian law on all inputs.

The observables $b^2,bh(x),h(x)^2$ are bounded by $B^2,B,1$.
Conditional Hoeffding, a union bound over $2|X|+1$ observables, and
successive conditioning over layers give all moment estimates with
probability at least $1-\kappa/2$ if
\[
 N\ge CB^4\alpha^{-2}\log[8H(2|X|+1)/\kappa].
\]
Stop controlling if $M\notin[1/2,3/2]$ and thereafter use the prior.
This defines the law on all histories, while the good event never stops.
Choose
\[
 \kappa=\min\{1/4,\tau/[32(A+1)]\},\qquad
 \alpha=\left[\frac{\kappa\sqrt{v_H}}{64He^K}\right]^2.
\]
Then the final coefficient has signed margin at least
$\sqrt{v_H}(1-\kappa/8)$. On this event shift the readout mean to
$ta$, where $t=(1+\kappa)/(2\sqrt{v_H})$; otherwise use zero mean.
Its mean output is at least $1/2+\kappa/4$ in the correct direction.
Its noise variance is at most $1/N$ at each input. If
$N\ge C\kappa^{-2}\log(4|X|/\kappa)$, the maximum noise is less
than $\kappa/8$ with probability at least $1-\kappa/2$. Thus the
controlled law $Q_N$ gives the strict target event probability at
least $1-\kappa$.

Use the same reference auxiliaries as in the program. The reference
weight marginal in~\eqref{ac-row-realization} remains Gaussian even
under the history-dependent orthogonal decomposition. Conditional
chain rule and data processing give hidden-block cost at most $ND_j$;
stopping only reduces this bound. The readout shift costs $Nt^2/2$.
After discarding auxiliaries,
\[
 \operatorname{KL}(Q_N\Vert\mu_N)
 \le N\left[\sum_jD_j+\frac{(1+\kappa)^2}{8v_H}\right]
 \le N(A+3A\kappa).
\]
Binary relative entropy for the target event implies
\[
 -N^{-1}\log\mu_N(E_f)
 \le\frac{A+3A\kappa+N^{-1}\log2}{1-\kappa}\le A+\tau
\]
when $N\ge8/\tau$. Since $v_H^{-1}\le8A$, the width bound
\eqref{ac-channel-width} implies all the numerical conditions above.
\end{proof}

\begin{theorem}[Critical raw target mass]\label{ac-target-mass}
For fixed $\phi$ satisfying the hypotheses of this appendix, there
are constants $C_0,C_1,C_2$ such that, if
$R=(d+L+2)A_F(f)$ and $L\ge C_0\log R$, then at every integer
$N\ge R^{C_1}$,
\[
 \mu_N\{f(x)z_W(x)>1/2\ (x\in X_d)\}
 \ge\exp\left[-N\left(V_*+C_2\frac{\log^2R}{L}\right)\right].
\]
\end{theorem}
\begin{proof}
Apply Lemma~\ref{ac-channel} to Lemma~\ref{ac-program} with
$\tau=1/H$. Each factor in~\eqref{ac-channel-width} is at most
$e^{CS}$: $B\le e^{CS}$, $K\le CS$, $H,d\le e^S$, and
$A+1\le C(1+S)$. Also $\log|X_d|=d\log2$. The additional
$1/H$ action is absorbed into $CS^2/H$. Enlarging constants handles
$H=L-1$ and $S=1+\log R$.
\end{proof}

\subsection{Near-baseline paths have parallel circuits}

For $D\ge0$ define $G(D)\ge1$ by
$G(D)-1-\log G(D)=2D$. Gaussian entropy duality implies
\begin{equation}\label{ac-gaussian-gain}
 \E_\rho\langle g,u\rangle^2\le G(D)\norm u^2,
 \qquad \sqrt{G(D)}\le e^{\sqrt D},\qquad G(D)\le4D+2
\end{equation}
for any tilt of a standard Gaussian with entropy at most $D$.
Indeed projection onto $\langle g,u\rangle/\norm u$ gives
$D\ge(m-1-\log m)/2$ when its second moment $m>1$; writing
$m=e^{2v}$ gives $(m-1-\log m)/2\ge v^2$.

\begin{lemma}[Variance and entropy budgets]\label{ac-occupation}
For a full-cube globally $1/2$-margined path of action
$A\le V_*+\epsilon$, $\epsilon\ge0$, put
$q_j=\max_xQ_j(x,x)$. Then
\[
 \sum_{j=0}^{H-1}q_j+\sum_{j=1}^H\sqrt{D_j}
          \le C[1+\log(H+2)+\sqrt{H\epsilon}].
\]
\end{lemma}
\begin{proof}
All one-input variances are positive, since a zero variance would
propagate at finite entropy and force zero final output. Decrease $c$
so $q-cq^2$ is increasing on $[0,2]$. Projecting onto an input attaining
the new maximum and using Lemma~\ref{ac-scalar} yields
\[
 D_j\ge q_j-q_{j-1}+cq_{j-1}^2.
\]
The scalar suffix at an input attaining $q_i$, including its readout,
costs at least $V(q_i)$. Thus
\begin{equation}\label{ac-prefix-budget}
 \sum_{j\le i}D_j\le\epsilon+Cq_i.
\end{equation}
These inequalities concern full layer entropies. No two-input
contraction property of the activation is needed.

For an interval of $t$ transitions whose pre-states and endpoint
are at least $u$, dividing dissipation by the squared pre-state and
using $(q'-q)/q^2\ge1/q-1/q'$ gives
$t\le C(E/u^2+1/u)$, where $E$ is its entropy expenditure.
An initial prefix of length $t$ and expenditure $E$ also obeys
\begin{equation}\label{ac-forced-cooling}
 \sum_{j<t}q_j\le C\log(t+2)+C\sqrt{tE}
\end{equation}
when total entropy is bounded by a fixed constant. To prove this,
write $F_0(v)=v-cv^2$, compare $q_j$ to $r_j+s_j$ with
$r_{j+1}=F_0(r_j)$, $s_{j+1}=F_0(s_j)+D_{j+1}$, $r_0=q_0\le2$,
$s_0=0$, and decrease $c$ so $F_0$ is increasing on the resulting
bounded range. Since $F_0(r+s)\le F_0(r)+F_0(s)$, induction proves
the comparison. The free sum is $O(\log(t+2))$, and
$c\sum s_j^2\le E$, proving~\eqref{ac-forced-cooling}.

If $\epsilon$ exceeds a fixed positive constant, summed dissipation
and Cauchy--Schwarz prove the claim. Otherwise total entropy is bounded.
Set $a=\max\{\epsilon,(H+1)^{-1}\}$, $b=\min_jq_j$, and
$u_0=\max\{a,b\}$. If $u_0$ exceeds a fixed small constant, the
occupation bound gives bounded $H$. Otherwise use dyadic levels
$u_j=2^ju_0$ until a fixed constant, and their last-exit indices
$\tau_j=\max\{i:q_i\le u_j\}$. There are $O(\log(H+2))$ bands.
The initial prefix has entropy $O(\epsilon+1/H)$ if $b\le a$;
then~\eqref{ac-forced-cooling} and Cauchy--Schwarz bound its two
contributions as claimed. If $b>a$, its entropy is $O(b)$ and its
length $O(1/b)$, giving state sum $O(\log(H+2))$ and square-root
entropy $O(1)$. A later nonempty last-exit block has entropy $O(u_j)$
by~\eqref{ac-prefix-budget}; after removing its first transition
its pre-states and endpoint exceed $u_j$, so its length is $O(1/u_j)$.
Every state in it is $O(u_j)$ because its starting state is at most
$u_j$ and any increase is at most spent entropy. Both contributions
are $O(1)$. The final tail has bounded length and contributions.
Summing proves the assertion.
\end{proof}

\begin{theorem}[Critical canonical circuit theorem]\label{ac-circuits}
A globally $1/2$-margined path of action $A\le V_*+\epsilon$ gives
a nonuniform threshold circuit of size $R_1^{O(1)}$ and depth
$O(\log R_1)$, and a bounded-fan-in Boolean circuit of size
$R_1^{O(1)}$ and depth $O(\log^2R_1)$, where
\[
 T=1+\log(H+2)+\sqrt{H\epsilon},\qquad
 R_1=(d+H+A+2)e^{CT}.
\]
All constants depend only on the fixed activation.
\end{theorem}
\begin{proof}
We give the anchor reduction, finite-dimensional state, evaluator,
and parallel iteration explicitly.

\paragraph{Anchor reduction.}
Let $I_j(U)$ be the entropy of the layer-$j$ preactivation marginal
on an anchor set $U$. It increases with $U$ and is at most $D_j$.
Choose $\rho=[C H(A+1)]^{-1}\le1/64$ and a cutoff $\delta>0$.
Greedily add a query whenever $I_j(U\cup\{x\})-I_j(U)>\delta$
for some layer, or projection of the final readout onto the anchor
feature span changes its output there by more than $\rho$.
There are at most $A/\delta$ entropy additions and $2r/\rho^2$
readout additions: each latter increases the projection's squared
norm by at least $\rho^2$, since final feature norms are at most one.

Replace each row density by its conditional expectation $w_j(G_U)$
given the anchors. It has entropy $I_j(U)\le D_j$ and preserves
the original anchor marginal exactly. The conditional query KL
increment is at most $\delta$. Conditional maximal coupling and
Pinsker give expected squared feature error at most
$4\sqrt{\delta/2}$, since features are bounded by one.
For kernels sharing their anchor block, compare a query by its
projection onto the common anchor span and its orthogonal residual
norm. This defines a metric $d_U$ equal to the RMS error in a Gaussian
coupling sharing the anchors and one residual standard Gaussian.
The preceding replacement has $d_U$ error $O(\delta^{1/4})$.
For a later Gaussian reference with the same anchor block, tilt this
one-query coupling by the common $w_j$. Its entropy is $I_j(U)$;
\eqref{ac-gaussian-gain} and the Lipschitz property multiply the
input error by at most $e^{\sqrt{D_j}}$. Thus the complete error is
at most $CH\delta^{1/4}\exp(\sum_j\sqrt{D_j})$.
Choose $\delta$ so this is at most
$\rho/\sqrt{\max\{2r,1\}}$. Lemma~\ref{ac-occupation} makes
the anchor count polynomial in $R_1$.

These replacements define a \emph{genuine} full-cube controlled path,
with the exact original anchor blocks, density depending only on
anchors, and no greater row entropies. The projected readout loses
at most $2\rho$ of margin. Scaling it by $(1-4\rho)^{-1}$ restores
margin $1/2$ at action increment $O(1/H)$. Its coefficient vector
has norm $O(\sqrt{A+1})$. The genuine path therefore satisfies
the budgets of Lemma~\ref{ac-occupation} with bound $CT$.

\paragraph{Canonical finite states.}
Initially $u_0(x)=(1,x)/\sqrt d\in\R^{d+1}$. Given states in
$\R^K$, let $g\sim\mathcal N(0,I_K)$ tilted by
$w_j((\langle g,u(a)\rangle)_{a\in U})$, and write this law as
$\rho_j$. In $L^2(\rho_j)$ put
$\Phi(u)=\phi(\langle g,u\rangle)$ and let $V$ span all $K$
coordinate functions $g_i$ and all anchor features $\Phi(u(a))$.
For an orthonormal basis $e_1,\ldots,e_M$, $M\le K+|U|$, define
\begin{equation}\label{ac-canonical}
 F(u)=\bigl(\langle e_b,\Phi(u)\rangle_{b=1}^M,
                         \norm{(I-P_V)\Phi(u)}\bigr).
\end{equation}
Pythagoras preserves every diagonal. At anchors the last coordinate
is zero, so every query-anchor cross moment and the entire anchor
block are preserved exactly. Future anchor-plus-one-query Gaussian
laws are determined by precisely these moments. Induction gives
$d_U=0$ relative to the genuine anchor-only path at every stage,
and identical anchored readout. The canonical full covariance is
not asserted to be a valid original transition. Its state dimension
is at most $d+1+H(|U|+1)$, polynomial in $R_1$.

Since the linear map $u\mapsto\langle g,u\rangle$ lies in $V$,
write $F(u)=A_ju+B_j(u)$, with
$\norm{A_j}_{\rm op}\le\alpha_j:=\sqrt{G(D_j)}$.
The map taking a Hilbert residual to its $V$ projection and its
orthogonal norm is $1$-Lipschitz. Moreover
$|\phi'(s)-1|\le C\min(s^2,1)$ by odd $C^3$ regularity and
$0\le\phi'\le1$. For $\norm u,\norm v\le\sqrt q$, put
$X=g\cdot u$, $Y=g\cdot v$, $Z=g\cdot(u-v)/\norm{u-v}$.
The squared normalized residual difference is at most a fixed
multiple of
$W=Z^2\min\{1,\max(|X|,|Y|)^4\}$.
For small fixed $c>0$,
\[
 e^{cW}-1\le cZ^2(X^4+Y^4)e^{cZ^2},\qquad
 \log\E_{\rm Gauss}e^{cW}\le Cq^2.
\]
For the second inequality condition $X$ or $Y$ on $Z$; its fourth
moment is at most $Cq^2(1+Z^4)$, regardless of correlations.
Entropy duality consequently gives
\begin{equation}\label{ac-residual-budget}
 \Lip(B_j)\le C(q+\sqrt{D_j}),\quad
 \sum_j\Lip(B_j)\le CT,\quad \prod_j\alpha_j\le e^{CT}
\end{equation}
on the actual state balls. Here use the genuine path's diagonal
budget and the original path's entropy budget; alternatively take
$D_j$ to be the actual anchor entropy. Extend each $B_j$ by radial
projection. Round the radius upwards on mesh $R_1^{-c_1}$ and impose
that same lower floor, with fixed large $c_1$. This adds only a
constant to the summed squared radii and preserves the actual path.

\paragraph{Polynomial finite evaluation.}
For any current law $\rho$, basis vector $e=(e_b)$, and polynomial
input radius, write
\[
 a(u)=\E_\rho e(g)\phi(g\cdot u),\quad
 b(u)=\E_\rho\phi(g\cdot u)^2,\quad
 F(u)=(a(u),\sqrt{b(u)-\norm{a(u)}^2}).
\]
Equation~\eqref{ac-gaussian-gain} gives
$\E_\rho\norm g^2\le K G(D)$. Truncate to $\norm g\le M_0$.
If the discarded probability is $p$, Bessel's inequality bounds the
discarded vector expectation uniformly by $\sqrt p$, and the scalar
expectation by $p$. A polynomial $M_0$ gives any inverse-polynomial
accuracy. On the retained set sample with importance density
$(1+\norm e^2)/Z_0$ relative to $\rho$, normalized there, where
$Z_0\le M+1$. The vector and scalar integrands have norms at most
$Z_0/2,Z_0$, and Lipschitz bounds at most $Z_0M_0/2,2Z_0M_0$.
An input net has logarithmic cardinality
$O(K\log(2+\poly(R_1)/\theta))$. For bounded Hilbert-valued
samples, expected mean error is at most the bound divided by the
square root of the sample count; bounded differences adds an
exponential tail. A union bound on this net therefore gives
polynomially many nodes approximating $a,b$ uniformly to error
$\theta$. Lipschitzness extends the estimate to the whole ball.
The retained nodes and all importance-weighted coefficients have
polynomial magnitude, irrespective of an anchor condition number.
Since $\norm a\le1$, choosing $\theta=O(\eta^2)$ and clipping
the approximate radicand at zero gives error $O(\eta)$ in $F$.

Round these constants, the matrix $A_j$, and all arithmetic to
inverse-polynomial accuracy. Scalar $\phi$ values need only a
polynomial mesh, since $\phi$ is $1$-Lipschitz; no algorithm for
evaluating $\phi$ is assumed. Every word has $O(\log R_1)$ bits.
Unary and binary word operations have polynomial-size truth tables.
A sum of polynomially many such words has polynomial integer range
after a common fixed-point scaling; threshold tests against each
possible integer and a bit decoder compute it in constant threshold
depth. Signed integer weights of polynomial magnitude can be
replaced by repeated inputs and complements. Squares, square roots,
comparisons, clipping, and radial projection have the same bounds;
division occurs only above the inverse-polynomial radius floor.
Thus one canonical residual evaluator has polynomial size and
constant threshold depth, uniformly on the required polynomial ball.

\paragraph{Parallel iteration.}
Set $S_j=\prod_{i\le j}\alpha_i$, $v_j=u_j/S_j$. The normalized
recursion is $v_j=C_jv_{j-1}+H_j(v_{j-1})$, where
$C_j=A_j/\alpha_j$ is a contraction and
$\Lip(H_j)\le\ell_j:=C(q_{j-1}+\sqrt{D_j})$ with
$\Lambda:=\sum_j\ell_j\le CT$. All $S_j$ are polynomial in $R_1$.
Precompute the interval products $C_{j:i}=C_j\cdots C_i$ as
nonuniform constant matrices; their operator norms and entries are
at most one. Starting with $v_j^{(0)}=C_{j:1}v_0$, compute all times
simultaneously by
\[
 v_j^{(k+1)}=C_{j:1}v_0+
               \sum_{i=1}^j C_{j:i+1}H_i(v_{i-1}^{(k)}).
\]
Strict time ordering gives
$\max_j\norm{v_j^{(k)}-v_j}\le2\sqrt2\Lambda^k/k!$:
expansion uses distinct ordered indices, whose products sum to at
most $\Lambda^k/k!$. Thus $O(\log R_1)$ iterations give any fixed
inverse-polynomial accuracy. Uniform additive error $\delta_0$ per
update is amplified by at most
$\sum_{k\ge0}\Lambda^k/k!=e^\Lambda=R_1^{O(1)}$.
The clipped residuals have bounded ranges, so all exact and inexact
iterates stay in a fixed polynomial ball. Rounded interval products,
quadrature, and arithmetic contribute such additive errors; divide
primitive tolerances by a polynomial to absorb the number of terms.

There are polynomially many coordinates and times, each update has
constant threshold depth, and there are $O(\log R_1)$ updates.
The final rescaling and readout amplify errors only polynomially;
choose precision so the output error is below $1/16$. This preserves
all signs. Replacing each polynomial-fan-in majority gate by a
logarithmic-depth Boolean summation circuit proves the Boolean bound.
\end{proof}

\subsection{A finite-event certificate and the full-cube ceiling}

\begin{lemma}[Bounded-dual extraction]\label{ac-dual}
On $m$ inputs put $k=m(m+1)/2$. For any raw event $E$ of critical
prior probability at least $e^{-NA}$, any $\tau>0$, and
$0<\Delta\le1$, there is a realization in $E$ and laws relative
to its empirical preceding Gaussian covariances such that
\[
 \max_{x,y}|\E_{\rho_j}\phi(G_x)\phi(G_y)-Q_j(x,y)|\le\Delta,
 \qquad \sum_jD_j+r\le A+\tau,
\]
provided
\begin{equation}\label{ac-dual-width}
 N\ge\frac C\tau\left[
 Hk\log\left(1+\frac{C(A/\tau+1)Hk}{\Delta}\right)
 +m\log(C(A/\tau+1))+1\right].
\end{equation}
Here $z$ is the realization's actual output and
$r=z^{\mathsf T}Q_H^\dagger z/2$. No symmetry of the event is needed.
\end{lemma}
\begin{proof}
For a reference feature law $\mu$ on $[-1,1]^m$, let $\Phi(h)$ list
its $k$ upper-triangular products and define
\[
 I_B(q)=\max_{\lambda\in[-B,B]^k}
       [\lambda\cdot q-\log\E_\mu e^{\lambda\cdot\Phi}].
\]
A maximizer exists. Its exponential tilt $\nu$ has moment vector
$a$; coordinatewise optimality on the box gives
$\lambda_i(q_i-a_i)=B|q_i-a_i|$. Hence
\begin{equation}\label{ac-dual-identity}
 I_B(q)=\operatorname{KL}(\nu\Vert\mu)+B\norm{q-a}_1\ge0.
\end{equation}
The same tilt of Gaussian preactivations has identical entropy and
moments. The objective is $2k$-Lipschitz in the sup norm of $\lambda$.
A mesh-$\eta$ grid has at most $(1+CB/\eta)^k$ points, and each
fixed exponential likelihood ratio for $N$ iid rows has mean one.
Thus for empirical moments $\widehat q$,
\[
 \E e^{N I_B(\widehat q)}\le e^{2Nk\eta}(1+CB/\eta)^k.
\]
This is uniform in the conditional reference. Iterate through layers.
Given the hidden layers, the actual readout covariance is $Q_H/N$,
so $\E[e^{(1-s)Nr}\mid\text{hidden}]\le s^{-m/2}$ for $0<s<1$.
Writing $I_j$ for the realized dual costs gives
\[
 \E e^{N[\sum I_j+(1-s)r]}
 \le e^{2NHk\eta}(1+CB/\eta)^{Hk}s^{-m/2}.
\]
Conditioning on $E$ has relative entropy at most $NA$. Entropy
duality bounds the conditional mean of $\sum I_j+(1-s)r$ by
\[
 A+2Hk\eta+N^{-1}[Hk\log(1+CB/\eta)+(m/2)\log(1/s)].
\]
Take $s=\tau/[16(A+\tau)]$, $\eta=\tau/(64Hk)$, and
$B=(A+\tau)/\Delta$. Under~\eqref{ac-dual-width}, some realization
in $E$ has weighted cost at most $A+\tau/8$, hence unweighted cost
at most $A+\tau$, since all $I_j\ge0$ and
$(A+\tau)(1-s)=A+15\tau/16$. Apply~\eqref{ac-dual-identity}
at each layer; its moment error is at most $I_j/B\le\Delta$.
\end{proof}

\begin{lemma}[Consistent-path repair]\label{ac-repair}
A certificate of total cost at most $B$ in Lemma~\ref{ac-dual}
has a genuine controlled path of cost at most $B$ and output
$\widehat z$ satisfying
\[
 \max_x|\widehat z(x)-z(x)|
            \le\sqrt{2B}\,H\sqrt{m\Delta}\,e^{\sqrt{HB}}.
\]
\end{lemma}
\begin{proof}
Two positive semidefinite $m$-point kernels differing entrywise by
at most $\Delta$ have square-root feature embeddings with row
distances at most $\sqrt{m\Delta}$. Their operator difference is
at most $m\Delta$, and
$\norm{K^{1/2}-{K'}^{1/2}}_{\rm op}
\le\sqrt{\norm{K-K'}_{\rm op}}$.
For the last inequality use $K\le K'+aI$, monotonicity of square
root, and $\sqrt{K'+aI}\le\sqrt{K'}+\sqrt a I$, then reverse
$K,K'$. Existing embeddings can be identified isometrically with
these square-root embeddings, adjoining orthogonal complements.

At each step couple Gaussian rows for the genuine preceding kernel
and the empirical reference in such a common feature space. Tilt
the joint Gaussian by the certificate density on its original
coordinate. The new marginal has no larger entropy. Equation
\eqref{ac-gaussian-gain} and the Lipschitz property multiply each
row RMS error by at most $e^{\sqrt{D_j}}$. Identify the controlled
feature moments with the next empirical kernel using the certificate
error, at additional row distance $\sqrt{m\Delta}$. Gluing the
shared feature spans gives
$e_j\le e^{\sqrt{D_j}}e_{j-1}+\sqrt{m\Delta}$, $e_0=0$.
Consequently $e_H\le H\sqrt{m\Delta}e^{\sqrt{HB}}$.
Transfer the empirical readout vector of norm $\sqrt{2r}$ in the
common space and project it onto the genuine final feature span.
Its norm does not increase, and its output changes by at most
$\sqrt{2r}e_H\le\sqrt{2B}e_H$.
\end{proof}

\begin{theorem}[Actual critical posterior circuit ceiling]\label{ac-ceiling}
Put $R_0=d+n+L+2$, and suppose $L\ge C\log R_0$. At the actual
critical scale $\sigma^2=1/N$, for every integer
$N\ge2^{Cd}R_0^C$ and every consistently labeled sample of at most
$n$ inputs, the original posterior gives probability at least
$31/32$ to a fixed full-cube class of functions having Boolean
circuits of size $\poly(R_0)$. More precisely this class can be taken
as $\{g:J_L(g)\le V_*+C\log^2R_0/L\}$; its threshold depth is
$O(\log R_0)$ and its Boolean depth is $O(\log^2R_0)$.
\end{theorem}
\begin{proof}
The memorizer $g_S=1-2\sum_{p\text{ distinct negative sample input}}
\boldsymbol1_{\{p\}}$ is Boolean, has the sample labels, and has
Fourier variation at most $1+2n$, since each point indicator has
Fourier variation one. Theorem~\ref{ac-target-mass} gives a raw
conditioning denominator at least $e^{-Na_0}$, where
$a_0=V_*+C_0\log^2R_0/L$, at all $N\ge R_0^{C_0}$ after
enlarging constants.

Let $E$ be the raw event of globally margined draws whose Boolean
function satisfies $J_L(g)>a_0+K_0/L$. This is measurable because
there are finitely many Boolean functions on the cube. If
$\mu_N(E)\ge e^{-N(a_0+1/L)}$, apply Lemma~\ref{ac-dual} on
the entire cube with $m=2^d$, action parameter $a_0+1/L$, and
$\tau=1/L$. Choose $\Delta$ so Lemma~\ref{ac-repair} has output
error at most $\rho=[C L(a_0+1)]^{-1}$. It suffices that
\[
 \log(1/\Delta)\le
 C[d+\log(L(a_0+2))+\sqrt{L(a_0+2)}]
\]
for an appropriate sufficiently small tolerance. The required width
in~\eqref{ac-dual-width} is therefore at most $2^{Cd}R_0^C$.
The repaired output keeps the selected signs; multiplying its readout
by $(1-2\rho)^{-1}$ restores margin $1/2$ and adds $O(1/L)$
action. This realizes a function in $E$ at action at most
$a_0+K_1/L$, contradicting $K_0>K_1$.

Hence $\mu_N(E)<e^{-N(a_0+1/L)}$. Division by the sample-consistent
denominator bounds posterior bad mass by $e^{-N/L}\le1/32$ under
the stated width threshold. Absorb $K_0/L$ into the logarithmic
excess. Theorem~\ref{ac-circuits} gives the claimed circuit bounds.
If the rate infimum is not attained, choose a path within $1/H$;
this changes the exponent $\sqrt{H\epsilon}$ only by a constant.
Every bound is simultaneous over all samples of the specified size,
and the class itself is sample independent.
\end{proof}


\section{A variant with every derivative globally bounded}
\label{app:all-derivatives}

The activation in Theorem~\ref{thm:main} has global bounds on any
prescribed finite number of derivatives. Its sufficiently high
derivatives need not be bounded. This feature is not necessary for an
exponential dependence on teacher depth: a smaller encoding amplitude gives
one activation whose derivatives of \emph{every} order are globally
bounded. Amplifier weights of order $d$ preserve the linear teacher depth.

\begin{theorem}\label{thm:all-derivatives}
Fix a finite integer $r\ge1$ and $\eta>0$. There is one computable
activation $\phi$ with all the shape and $C^r$-closeness properties of
Theorem~\ref{thm:main}, such that every derivative of $\phi$ is globally
bounded. The same circuit-hard targets $f_d$ have globally unit-margin
teachers with
\[
 L_d=d+4,\qquad
 \text{active width }O(\log(d+2)),\qquad
 \text{coefficient magnitude }O(d+1),
\]
and polynomial coefficient bit length. For some $c>0$, every sample
of at most $\exp(cd)$ observations has posterior confidence-failure
rate at least $139/1024$, at every fixed Gaussian exponent
$\gamma\ge1$ and every sufficiently large width
$N\ge\exp(C_\gamma d)$. The fixed-width cold-limit conclusion holds
at every width above the memorizer threshold~\eqref{eq:MB}.

Thus this variant retains the sample lower bound
$\exp(\Omega(L_d))$. It does not assert a uniform
analytic-strip bound, a finite-order Gevrey bound, or efficient
activation evaluation.
\end{theorem}

\subsection{The modified bands and derivative estimates}

Fix a sufficiently large integer $M_0\ge24$, and set
\begin{align*}
 m_d&=M_0+\lceil\log_2(d+2)\rceil,&
 J_d&=2^{m_d+10},&
 \Delta_d&=2^{-d-4},\\
 r_d&=\Delta_d/J_d,&
 c_{d,k}&=-4d-\tfrac14-\Delta_d k,&
 \delta_d&=\tau r_d2^{-(m_d-3)d}.
\end{align*}
In particular,
\begin{equation}\label{eq:all-derivative-amplitude}
 \delta_d=\tau2^{-(m_d-2)d-m_d-14}.
\end{equation}
Use these parameters in~\eqref{eq:S}--\eqref{eq:phi}, with a fixed
sufficiently small positive rational $\tau$. The truth tables and
centers are otherwise unchanged.

For each fixed derivative order $j$, the Gaussian lattice estimate
in Lemma~\ref{lem:phi} gives
\begin{equation}\label{eq:all-derivative-bound}
 \sup_{\R}|S_d^{(j)}|
 \le C_j\tau(16J_d)^{j-1}2^{-(m_d-j-2)d}.
\end{equation}
Here $S_d$ denotes the whole $d$th band. This is summable in $d$
for every fixed $j$, because $m_d=\log_2d+O(1)$ and $J_d=O(d)$.
Consequently every derivative is globally bounded. On each complex
compact, the moving band centers and Gaussian suppression still
give locally uniform convergence; hence $S$ is entire. The explicit
tails and the computable finite truth tables give computability.

The weighted first- and second-derivative estimates have band
majorants
\[
 C\tau e^{8d}2^{-(m_d-3)d},\qquad
 C\tau J_de^{8d}2^{-(m_d-4)d},
\]
respectively. Both are summable. The cubic factor at zero supplies
the same extra powers of $u$ as in Lemma~\ref{lem:phi}. Therefore
\eqref{eq:weightedfirst}--\eqref{eq:weightedsecond} hold, with fixed
constants for this modified activation. A sufficiently small $\tau$
preserves oddness, slope one at zero, strict increase, global
Lipschitz constant one, positive-half-line strict concavity, and
saturation to $\pm1$. Equation~\eqref{eq:all-derivative-bound} makes
the prescribed $C^r$ distance smaller than $\eta$ after a further
fixed reduction of $\tau$. The constants bounding derivative $j$
can grow rapidly with $j$; no uniform analytic-type estimate is used.

\subsection{A growing finite-difference decoder}

Use $h_d=\Delta_d/(2m_d)$ and the $m_d$th finite difference in
\eqref{eq:decoder}. At $c_{d,k}$ the central bump contributes
$\delta_df_d(k)$. Every other query is at least $h_d$ from every
center in this band. Since $J_d=2^{m_d+10}$, its leakage, multiplied
by the coefficient sum $2^{m_d}=O(d)$, is $o(\delta_d)$.

For another band $e\ne d$, the logarithm of its amplitude ratio
is $O(d\log d+e\log e)$, whereas Gaussian suppression has exponent
at most $-c(d-e)^24^eJ_e^2$. Split the sum into $e\le d/2$,
$d/2<e<d$, and $e>d$. In the first range even the smallest fixed
$e$ gives a negative quadratic in $d$, dominating $d\log d$.
The second has suppression at least $\exp(-c4^{d/2})$; in the
third the amplitudes decrease and the Gaussian suppression is
stronger still. Polynomial factors and the coefficient sum do not
change these conclusions. Reflected bands are farther away.

For $m=m_d=O(\log d)$, the negative-half-line exponential series of
tanh gives
\[
 \sup_{u\le-4d+1}|\tanh^{(m)}u|
 \le C2^m e^{-8d+2}.
\]
Indeed the differentiated series is bounded by
$2^{m+1}\sum_{n\ge1}n^me^{-(8d-2)n}$, whose first term dominates
uniformly for large $d$ since $m=o(d)$. The tanh contribution to
the finite difference is at most
$C\Delta_d^m m^{-m}e^{-8d}$. Dividing by
\eqref{eq:all-derivative-amplitude} bounds its relative size by
\[
 C\tau^{-1}2^{-2d-3m+14}m^{-m}e^{-8d}=o(1).
\]
We therefore retain the uniform decoding estimate
$F_d(c_{d,k})=\delta_d(f_d(k)+o(1))$.

The first layer has $m_d+1$ neurons and evaluates these query
points. The next layer uses binomial coefficients of magnitude at
most $2^{m_d}=O(d)$. Set the amplifier gain
\[
 a_d=2^{m_d-2}=\Theta(d).
\]
Exactly $d+1$ scalar amplification layers give ideal final amplitude
\[
 \delta_da_d^{d+1}=\tau2^{-16}=: \rho.
\]
Their ideal preactivation amplitudes form a geometric sequence with
sum at most $2\rho$, since $a_d\ge2$. The estimate
\eqref{eq:taylor} therefore gives a uniform $O(\rho)$ relative-error
bootstrap, independent of the growing gain. Choosing $\tau$ small
keeps every actual preactivation below $2\rho$; a fixed readout
$2/\rho$ gives margin greater than one. Counting the query layer,
combination layer, $d+1$ amplifiers, and readout gives $L_d=d+4$.
The amplifier weights are still $O(d)$. This proves the teacher
resource bounds. Rounding its coefficients to $O(d\log d)$ bits
preserves the margin: the whole propagation gain is at most
$\exp(O(d\log d))$.

\subsection{Cold and fixed-gap posterior bounds}

Let $a>0$ be the hard-product exponent in the proof of
Lemma~\ref{lem:hard}. Since $L_d=d+4$, the original fixed-gain
memorizer applies. Set
\[
 q=\exp(a/100),\qquad \alpha=c_0(q-1),
\]
with the same sufficiently small constant $c_0$ as in
Lemma~\ref{lem:memorizer}. Both are fixed, and
$q^{L_d}=\exp(ad/100+O(1))$.
The relative-error sum is still $O(\alpha/(q-1))$, hence
uniformly small. The first scalar energy is at most one: its
decaying factor is
\[
 q^{-2(L_d-2)}=\exp(-ad/50+O(1)),
\]
whereas \eqref{eq:B0} gives $\log B_0\le3cd+O(\log d)$.
Choose $c<a/1000$. The witness energy satisfies
\begin{equation}\label{eq:all-derivative-memorizer}
 K_d\le q^2L_d+O(1).
\end{equation}
The cold posterior argument remains valid at each fixed
architecture. Arithmetic--geometric mean bounds its unit-margin
representative product by
\[
 \exp\{ad/100+O(1)\}\le e^{ad}
\]
for all large $d$. This is the same hard-table threshold as before.

For the actual fixed-gap scales, Lemma~\ref{lem:perturb} and
Proposition~\ref{prop:finite} apply with the energy in
\eqref{eq:all-derivative-memorizer}. Since $L_d=d+4$ and
$q^{L_d}=\exp(ad/100+O(1))$, the inverse perturbation tolerance
remains $\exp(O(d))$ without further changes. The Gaussian bulk
constant is polynomial in $d,L_d$, and the memorizer width
in~\eqref{eq:MB} is $\exp(O(d))$ uniformly over $n\le e^{cd}$.
Thus every $N\ge\exp(C_\zeta d)$ meets the finite-scale
condition~\eqref{eq:width}, for each fixed gap $\zeta=\gamma-1>0$.
The posterior product bound remains below $e^{ad}$; the hard-table
and Gibbs-to-confidence calculation gives the same failure rate
$139/1024$.

\subsection{The critical endpoint}

The modified activation satisfies the fixed-activation hypotheses
of Appendix~\ref{ac-appendix}: the same shape and weighted estimates,
bounded derivatives through order five, and the strict scalar
potential obtained by choosing $\tau$ sufficiently small. Therefore
its finite-width target-mass and canonical-circuit results apply
with constants depending only on this fixed activation.

For $L_d=d+4$ and $n\le e^{cd}$, reducing $c$ makes the depth
condition $L_d\ge C\log(d+n+L_d)$ hold for all large $d$.
A further reduction makes the posterior circuit ceiling
$\poly(d+n+L_d)$ at most
$2^{d/4}$. The width threshold
$2^{Cd}(d+n+L_d)^C$ is still at most $\exp(C'd)$.
The hard-table and confidence calculation is unchanged. This
completes the proof of Theorem~\ref{thm:all-derivatives}.

\subsection{Quantitative derivative bounds and numerical evaluation}

The distinction between bounded derivatives of each order and a
quantitative bound on their growth is important here.

\begin{proposition}\label{prop:gevrey-evaluation}
Suppose a real activation satisfies, for fixed constants $C,A\ge1$
and finite $s\ge1$,
\[
 \sup_{u\in\R}|\phi^{(j)}(u)|\le CA^j(j!)^s
 \qquad(j\ge0).
\]
On dyadic inputs of bit length $b$ in $[-R,R]$, it has nonuniform
Boolean circuits approximating its value to error $2^{-p}$, of
size polynomial in $R+1$, $p+1$, and $b+1$. The polynomial may
depend on $C,A,s$. No efficient uniform procedure for computing
the circuit's coefficients is assumed.
\end{proposition}
\begin{proof}
Replace $s$ by the integer $s_0=\lceil s\rceil$, and choose an
integer upper bound for $A$. For an integer $n$, partition the
interval into pieces of radius at most
$h=(2An^{s_0-1})^{-1}$. Taylor expansion of degree $n-1$ at the
center of a piece has remainder at most
\[
 \frac{CA^n(n!)^{s_0}}{n!}h^n
 \le C(Ahn^{s_0-1})^n=C2^{-n}.
\]
Take $n=O(p+\log C+1)$, with enough fixed slack. There are
$O((R+1)n^{s_0-1})$ pieces. Store their centers and Taylor
coefficients as nonuniform circuit constants. The coefficients
have logarithmic magnitude $O(n\log(n+1))$, and the local
argument has magnitude at most one. Rounding each coefficient
to absolute error $2^{-p}/(8n)$ controls the total coefficient
error. Polynomially many additional bits suffice for the dyadic
arithmetic, interval selection, and evaluation of each polynomial.
All sizes are polynomial in the displayed resources.
\end{proof}

This includes a global finite-order Gevrey bound. A bounded
holomorphic extension to a fixed-width strip is another sufficient
condition, by Cauchy estimates. Such conditions cannot hold for
the constructed activation: simulating its polynomial-resource
teacher at precision $O(d\log d)$ would give polynomial-size
Boolean circuits for the circuit-hard $f_d$. Thus the new variant
still contains computationally expensive information in its
activation primitive. It is not a counterexample for a standard
efficiently evaluable activation or for exact tanh.


\end{document}
